% ---------------------------------------------------------------------------
% Artifact citation block for the CoreGate paper.
%
% CoreGate: Execution-Lineage Supervision for Attack Technique Recognition
% under Environment Shift
% Anonymous Authors (double-blind review copy)
%
% How to use this file
%   1. Paste the two BibTeX entries below into paper/coregate_refs.bib, or into
%      the camera-ready bibliography. They are camera-ready material: they name
%      the GitHub owner and the Zenodo DOIs, so they must not appear in the
%      review PDF.
%   2. Replace the placeholders once they exist. In this review copy the owner is
%      already the review account and Option A already carries the registered
%      review link (/r/coregate-6342/), so only the DOI placeholders remain:
%        XXXXXXX/YYYYYYY   -> the software / dataset Zenodo DOIs
%   3. Paste exactly one data-availability paragraph into the paper: Option A
%      while the submission is under double-blind review, Option B for the
%      camera-ready.
%   4. Cite the software with \cite{coregate2027artifact}. Do not cite the paper
%      and the artifact with the same key.
%
% Option A is the only paragraph that may appear in the review PDF: it names no
% author, no institution and no repository owner. Option A needs \usepackage{url}
% (or hyperref) for \url{...}.
%
% The dataset DOI resolves only after acceptance: the paper states that code is
% released publicly while the attack-emulation dataset follows acceptance.
% ---------------------------------------------------------------------------

% ---- paste into paper/coregate_refs.bib (camera-ready only) ---------------

@software{coregate2027artifact,
  title     = {{CoreGate}: Execution-Lineage Supervision for Attack Technique
               Recognition under Environment Shift},
  author    = {{Anonymous Authors}},
  year      = {2027},
  version   = {0.7},
  publisher = {Zenodo},
  doi       = {10.5281/zenodo.XXXXXXX},
  url       = {https://github.com/holdenli275/coregate}
}

@misc{coregate2027dataset,
  title     = {{CoreGate} {Linux} Telemetry Candidate Set (template v2)},
  author    = {{Anonymous Authors}},
  year      = {2027},
  version   = {2.0},
  publisher = {Zenodo},
  doi       = {10.5281/zenodo.YYYYYYY},
  note      = {360 labeled attack-emulation graphs over 45 {ATT\&CK}
               sub-techniques plus 14 unlabeled entries; collected by the
               authors by executing Atomic Red Team tests in a controlled
               Linux lab. Released upon acceptance.}
}

% ---- paste into the paper -------------------------------------------------

% Option A - double-blind review version. Use this while the submission is
% anonymous. The link below is the registered review mirror; do not add the
% GitHub owner, the Zenodo DOIs, or the artifact BibTeX entry to the review PDF,
% because each of them identifies the authors.
% Code and data availability. The implementation, the dataset catalog and the
% released per-seed experiment outputs are available to reviewers at
% \url{https://anonymous.4open.science/r/coregate-6342/}. The attack-emulation
% dataset is a separate deposit and will be released publicly, together with the
% code, upon acceptance.

% Option B - camera-ready version. Replace the placeholders first. Every number
% in the paper is reproducible from the released per-seed run JSONs without the
% dataset.
% Code and data availability. The implementation and the released experiment
% outputs are publicly available at~\cite{coregate2027artifact}. The
% attack-emulation dataset is a separate deposit and will be released upon
% acceptance. Evidence boundaries are automatically extracted weak labels from
% the collection journals, not independent manual annotation, so the catalog
% documents a within-catalog recognition gain rather than a formal benchmark.

% Option C - the sentence the submitted manuscript currently carries. It is
% already author-free and stays valid if the venue forbids external links
% during review.
% Code and data availability. Code will be released publicly; the
% attack-emulation dataset will follow upon acceptance.
